\section{对联生成样例}\label{app:couplets}

下表选取14个评测上联，按联长排列，并列数据集参考下联与两种解码结果。两组输出来自相同的最终微调模型和双向评分；区别在于是否启用外部规则。参考下联为配对语料中的一个答案，生成文本保留原始用字。

\begingroup
\small\setlength{\tabcolsep}{4pt}\renewcommand{\arraystretch}{1.1}
\noindent\begin{tabularx}{\linewidth}{@{}C{3.65cm}C{3.65cm}YY@{}}
\toprule
\textbf{上联}&\textbf{参考下联}&\centering\arraybackslash\shortstack{\textbf{微调模型}\\无规则解码}&\centering\arraybackslash\shortstack{\textbf{微调模型}\\规则约束解码}\\
\midrule
\rule{0pt}{9.4mm}鸿运连交&礼花才放&祥光普照&祥光普照\\
\rule{0pt}{9.4mm}春归追梦处&福至赏梅时&马到到头时&秋至落花时\\
\rule{0pt}{9.4mm}且舒敞胸襟&应放开眼界&且放开眼界&更放开眼界\\
\rule{0pt}{9.4mm}诗歌伴雨吟&联墨随绳舞&墨舞随风舞&画卷随风舞\\
\rule{0pt}{9.4mm}山好只因人&水灵曾有幸&月圆只为客&水清还为客\\
\rule{0pt}{9.4mm}有愿即参禅&通情方达理&无心即悟道&无心方悟道\\
\rule{0pt}{9.4mm}爽气满胸襟&清风生掌握&清风盈袖袖&清风盈眼界\\
\rule{0pt}{9.4mm}一园春烂漫&八秩史辉煌&百业业辉煌&万里业辉煌\\
\rule{0pt}{9.4mm}经春秋冬夏&赴南北西东&历春夏秋冬&历南北东西\\
\rule{0pt}{9.4mm}看砥柱中流&听鼓角午夜&喜神州大地&喜神州大地\\
\rule{0pt}{9.4mm}夜宴方阑归碧落&夕阳虽好近黄昏&春风又到到黄昏&春风正好到黄昏\\
\rule{0pt}{9.4mm}菊梦依稀秋意远&莲心自在月光清&枫情依旧月光寒&梅魂寂寞月光寒\\
\rule{0pt}{9.4mm}旧俗已随旧岁去&新生应逐新年来&新春更伴新春来&新年更待新春来\\
\rule{0pt}{9.4mm}无花无酒无花酒&有洞有天有洞天&有水有山有水人&有雨有风有雨风\\
\bottomrule
\end{tabularx}\par
\endgroup

\vspace{9pt}
\noindent\textbf{样例观察。}“春归追梦处”“爽气满胸襟”呈现重复字的变化；“菊梦依稀秋意远”呈现同位用字的变化。“鸿运连交”在两种解码下均为“祥光普照”，显示规则也可保留原候选。“无花无酒无花酒”的重复模式得到保留，内容组织及词语搭配的质量则由具体文本呈现。
